\section{Paire de lecteurs généalogiques sur l'atlas complet}
\label{sec:operadores-masa-two-way-rev11}
\index{masse!opérateurs bidirectionnels}
\index{atlas!lecteurs chronologique et dodécaphasique}

Le parcours enrichi possède deux représentations internes complémentaires,
\[
 \Gamma\longmapsto
 \bigl(\gamma_{108}(\Gamma),\tau_{12}(\Gamma)\bigr)
 \in\mathbb F_3^{108}\times\{-1,0,+1\}^{12}.
\]
La première conserve l'histoire tritique complète ; la seconde conserve
l'enregistrement dodécaphasique signé. Chaque projection induit un lecteur naturel, et les deux
sont évalués avant l'introduction des masses, des étiquettes PDG ou des résidus.

\HMTClaim{MD-R-MASS-TWOWAY-READER-PAIR}
\subsection{Paire de lecteurs généalogiques \(\mathbf K(\Gamma)=(K_D(\Gamma),K_\Omega(\Gamma))\) : domaine, codomaine et évaluations scalaires}

Avec la notation des
\cref{eq:KD-rev11,eq:kappaD-rev11,eq:kappa-omega-rev11}, nous définissons
\begin{equation}
 \mathbf K(\Gamma)
 :=\bigl(K_D(\Gamma),K_\Omega(\Gamma)\bigr)
 \in\left(\mathbb Z^2\times\tfrac12\mathbb Z\right)
 \times\mathbb Z^3,
 \label{eq:lector-conjunto-masas-rev11}
\end{equation}
où
\[
 K_D=\left(D_{27}+D_{36},D_1,\frac{D_4-D_3}{2}\right),
 \qquad
 K_\Omega=(\Omega_{90\mid120},54,P_6).
\]
Les évaluations scalaires sont
\begin{align}
 \kappa_D(\Gamma)
 &=D_{27}+D_{36}-\alpha_{\rm HMT}D_1
 +\frac{\Delta_4}{2}(D_4-D_3),
 \label{eq:kappaD-general-rev11}\\
\kappa_\Omega(\Gamma)
&=\Omega_{90\mid120}(\tau_\Gamma)
-54\alpha_{\rm HMT}+P_6(\tau_\Gamma)\Delta_4.
\label{eq:kappaOmega-general-rev11}
\end{align}
La deuxième coordonnée exige une distinction causale. Dans APP, le saut
spectral primordial est \(9-5=4\) ; le coefficient local \(2\), la compression
\(6=51-45\) et le déploiement bidimensionnel \(54=9\cdot6=459-405\) sont
des invariants ultérieurs et typés, conformément au
\cref{thm:app-jerarquia-cuatro-dos-seis-cincuentaycuatro}. Dans le lecteur
chronologique des déplacements, en revanche, la coordonnée électronique est
\(B_D=D_1(\Gamma_e)=54\) : elle compte les frontières immédiates et coïncide avec la
périodicité cinématique d'ordre \(54\) de l'endomorphisme CT108. Le
\(54\) fixe de
\(K_\Omega\) est la coordonnée commune transférée lorsque l'on rend compatibles les deux
lecteurs dans la réduction électronique. Son égalité numérique avec le défaut
global d'APP constitue un contrôle croisé ultérieur ; elle ne fait pas du
nombre \(54\) la différence primordiale entre somme et produit.

L'égalité \(4\cdot90=3\cdot120=360\) explique la généalogie commune des
deux lecteurs. \(K_D\) compare les déplacements du mot tritique
\(\gamma_{108}\) ;
\(K_\Omega\) compare les énergies de fenêtres dans le mot de douze phases. La
coïncidence d'origine n'identifie pas leurs quotients d'information.
Dans le domaine certifié \(\Gamma_{324}^{\rm or}\), la périodicité de
demi-tour rend pairs tous les \(D_k\), de sorte que
\((D_4-D_3)/2\in\mathbb Z\) et \(K_D\in\mathbb Z^3\). Hors de ce domaine,
la troisième composante conserve le type \(\tfrac12\mathbb Z\).

\subsection{Naturalité des opérateurs bidirectionnels et préservation des types}

Le lecteur \(D_k\) est invariant par translation temporelle, renversement et toute
permutation bijective de l'alphabet tritique. Dans les parcours déclarés,
l'inversion simultanée des deux orientations agit par une translation de
27 pas discrets et conserve \(K_D\). Le lecteur \(K_\Omega\) est invariant par
rotation et renversement dodécaphasiques, ainsi que par
\(C_M(\tau)=-\operatorname{rev}(\tau)\). Sa forme quadratique
\[
 \Omega(\tau)=\langle\tau,
 (4W_3^*W_3-3W_4^*W_4)\tau\rangle
\]
possède une polarisation bilinéaire entière et satisfait l'identité de 2-cocycle.

Les compteurs globaux de rotation, de changement de feuillet, de répétition de phase, de changement de
préfixe et de retour sont constants sur les 324 parcours. Par conséquent, toute
fonctionnelle qui se factorise exclusivement par ces six compteurs est constante
et se simplifie dans les rapports. Les lecteurs de
\eqref{eq:kappaD-general-rev11} et de
\eqref{eq:kappaOmega-general-rev11} agissent sur l'histoire locale et satisfont
ce contrôle de non-factorisation.

\begin{ablation}[Nécessité des deux feuillets dans le parcours électronique]
\label{abl:hojas-lector-two-way-rev11}
Le lecteur chronologique produit
\[
 (A_D,B_D,C_D)=
 \begin{cases}
 (80,54,6),&\substack{\text{mot chronologique conjoint des deux feuillets ;}\\
 B_D=D_1(\Gamma_e)=54},\\
 (80,84,-8),&\text{projection additive},\\
 (80,60,13),&\text{projection multiplicative}.
 \end{cases}
\]
Le triplet électronique complet exige l'histoire conjointe des deux feuillets
arithmétiques.
\end{ablation}

\subsection{Dénombrements exacts sur les 324 parcours}

Le lecteur chronologique produit quatorze profils. Le tableau suivant en donne le
dénombrement complet ; \(A_D=D_{27}+D_{36}\), \(B_D=D_1\) et
\(C_D=(D_4-D_3)/2\).

\begin{longtable}{@{}rrr r@{}}
\caption{Dénombrement exact des quatorze profils du lecteur chronologique sur les
324 parcours orientés.}
\label{tab:censo-perfiles-kd-rev11}\\
\toprule
\(A_D\)&\(B_D\)&\(C_D\)&multiplicité\\
\midrule
\endfirsthead
\toprule
\(A_D\)&\(B_D\)&\(C_D\)&multiplicité\\
\midrule
\endhead
0&24&0&18\\
40&24&2&36\\
40&54&6&18\\
40&78&27&36\\
40&84&30&36\\
60&78&25&18\\
60&84&28&18\\
80&54&6&18\\
80&54&7&18\\
80&66&2&18\\
80&66&3&36\\
80&66&4&18\\
100&66&0&18\\
100&66&2&18\\
\midrule
\multicolumn{3}{r}{total}&324\\
\bottomrule
\end{longtable}

Le lecteur dodécaphasique produit neuf profils :

\begin{table}[H]
\centering
\caption{Dénombrement exact des neuf profils du lecteur dodécaphasique sur
les 324 parcours orientés.}
\label{tab:censo-perfiles-komega-rev11}
\begin{tabular}{@{}rrr r@{}}
\toprule
\(\Omega\)&\(B_\Omega\)&\(P_6\)&multiplicité\\
\midrule
80&54&6&198\\
56&54&5&68\\
48&54&5&34\\
36&54&4&8\\
20&54&4&4\\
4&54&2&4\\
40&54&4&4\\
24&54&4&2\\
24&54&2&2\\
\midrule
\multicolumn{3}{r}{total}&324\\
\bottomrule
\end{tabular}
\end{table}

Les lecteurs coïncident comme triplets sur (18) parcours, toujours avec la valeur
\((80,54,6)\), et la paire \(\mathbf K\) prend (40) valeurs. Dans la section
canonique \(B1\) à 81 positions, \(K_D\) est constant dans 41 des 56
familles et raffine le quotient en 77 classes ; \(K_\Omega\) est constant dans 50
et produit 62 classes ; la paire en produit 78. Dans le relèvement à 324
orientations, \(K_D\) produit 146 classes et n'est constant dans aucune des
familles ; \(K_\Omega\) en produit 100 et descend dans 17 familles ; la paire en produit
151 et n'est constante dans aucune. Ces dénombrements prouvent que les deux
projections retiennent des mémoires distinctes et fixent le domaine de chaque
affirmation de descente.

\subsection{Traces vectorielles et spectres conjoints des \(13\) multisections}

Dans la section canonique \(B1\), nous notons
\(\ell_j,\nu_j,d_j,u_j\) les multisections de génération \(j\) des
secteurs leptonique chargé, neutrino, quark de type down et quark de type
up, respectivement. Pour chaque lecteur
\(r\in\{D,\Omega\}\) et chaque composante \(a\in\{1,2,3\}\), nous définissons
l'opérateur diagonal
\[
 K_{M,a}^r e_\Gamma=K_{r,a}(\Gamma)e_\Gamma,
 \qquad
 \mathbf K_M^r=(K_{M,1}^r,K_{M,2}^r,K_{M,3}^r).
\]
Les trois opérateurs de \(\mathbf K_M^r\) sont autoadjoints et commutent. Leur
trace normalisée s'entend donc composante par composante,
\[
 \overline{\mathbf K}_M^r
 =\frac{1}{\dim M}
 \bigl(\operatorname{Tr}K_{M,1}^r,
       \operatorname{Tr}K_{M,2}^r,
       \operatorname{Tr}K_{M,3}^r\bigr),
\]
et leur spectre est le spectre conjoint, c'est-à-dire la distribution avec
multiplicité des triplets simultanés. Les treize traces vectorielles sont :

\begin{longtable}{@{}c r c c@{}}
\caption{Traces vectorielles normalisées des lecteurs \(K_D\) et
\(K_\Omega\) dans les treize multisections canoniques.}
\label{tab:trazas-multisecciones-two-way-rev11}\\
\toprule
\(M\)&\(\dim M\)&\(\overline{\mathbf K}_M^D\)&\(\overline{\mathbf K}_M^\Omega\)\\
\midrule
\endfirsthead
\toprule
\(M\)&\(\dim M\)&\(\overline{\mathbf K}_M^D\)&\(\overline{\mathbf K}_M^\Omega\)\\
\midrule
\endhead
\(\ell_0\)&1&\((80,54,6)\)&\((80,54,6)\)\\
\(\ell_2\)&8&\((65,249/4,17/2)\)&\((80,54,6)\)\\
\(\ell_4\)&16&\((225/4,489/8,21/2)\)&\((141/2,54,11/2)\)\\
\(\nu_1\)&2&\((40,51,29/2)\)&\((60,54,5)\)\\
\(\nu_2\)&4&\((45,63,29/2)\)&\((61,54,5)\)\\
\(\nu_3\)&6&\((170/3,55,34/3)\)&\((206/3,54,11/2)\)\\
\(\nu_4\)&8&\((105/2,237/4,95/8)\)&\((143/2,54,45/8)\)\\
\(d_1\)&2&\((40,51,29/2)\)&\((60,54,5)\)\\
\(d_2\)&4&\((45,63,57/4)\)&\((61,54,5)\)\\
\(d_3\)&6&\((170/3,55,23/2)\)&\((218/3,54,17/3)\)\\
\(d_4\)&8&\((105/2,237/4,49/4)\)&\((149/2,54,23/4)\)\\
\(u_1\)&4&\((65,66,9)\)&\((80,54,6)\)\\
\(u_3\)&12&\((205/3,119/2,113/12)\)&\((74,54,23/4)\)\\
\bottomrule
\end{longtable}

La notation \((a,b,c)^{\times n}\) indique la multiplicité \(n\). Le spectre
conjoint complet de \(\mathbf K_D\) est :

\begingroup\footnotesize
\begin{longtable}{@{}c p{.82\textwidth}@{}}
\caption{Spectres conjoints complets du lecteur \(K_D\) dans les treize
multisections canoniques.}
\\
\toprule
\(M\)&\(\operatorname{Spec}_{\rm conj}\mathbf K_M^D\)\\
\midrule
\endhead
\(\ell_0\)&\((80,54,6)^{\times1}\)\\
\(\ell_2\)&\((0,24,0)^{\times1};(40,78,27)^{\times2};(80,54,7)^{\times1};(80,66,2)^{\times1};(80,66,3)^{\times1};(100,66,0)^{\times1};(100,66,2)^{\times1}\)\\
\(\ell_4\)&\((0,24,0)^{\times1}\); \((40,24,2)^{\times2}\); \((40,54,6)^{\times2}\); \((40,84,30)^{\times2}\); \((60,78,25)^{\times3}\); \((80,66,2)^{\times2}\); \((80,66,3)^{\times3}\); \((80,66,4)^{\times1}\)\\
\(\nu_1\)&\((40,24,2)^{\times1};(40,78,27)^{\times1}\)\\
\(\nu_2\)&\((0,24,0)^{\times1};(40,84,30)^{\times1};(60,78,25)^{\times1};(80,66,3)^{\times1}\)\\
\(\nu_3\)&\((40,24,2)^{\times2};(40,78,27)^{\times1};(60,84,28)^{\times1};(80,54,6)^{\times1};(80,66,3)^{\times1}\)\\
\(\nu_4\)&\((0,24,0)^{\times1};(40,24,2)^{\times1};(40,78,27)^{\times2};(40,84,30)^{\times1};(80,54,7)^{\times1};(80,66,2)^{\times1};(100,66,0)^{\times1}\)\\
\(d_1\)&\((40,24,2)^{\times1};(40,78,27)^{\times1}\)\\
\(d_2\)&\((0,24,0)^{\times1};(40,84,30)^{\times1};(60,78,25)^{\times1};(80,66,2)^{\times1}\)\\
\(d_3\)&\((40,24,2)^{\times2};(40,78,27)^{\times1};(60,84,28)^{\times1};(80,54,6)^{\times1};(80,66,4)^{\times1}\)\\
\(d_4\)&\((0,24,0)^{\times1};(40,24,2)^{\times1};(40,78,27)^{\times2};(40,84,30)^{\times1};(80,54,7)^{\times1};(80,66,3)^{\times1};(100,66,2)^{\times1}\)\\
\(u_1\)&\((40,54,6)^{\times1};(60,78,25)^{\times1};(80,66,2)^{\times1};(80,66,3)^{\times1}\)\\
\(u_3\)&\((40,24,2)^{\times2}\); \((40,78,27)^{\times2}\); \((60,84,28)^{\times1}\); \((80,54,6)^{\times3}\); \((80,66,3)^{\times1}\); \((80,66,4)^{\times1}\); \((100,66,0)^{\times1}\); \((100,66,2)^{\times1}\)\\
\bottomrule
\end{longtable}
\label{tab:espectros-kd-multisecciones-rev11}
\endgroup

Le spectre conjoint complet de \(\mathbf K_\Omega\) sur les mêmes fibres est :

\begingroup\footnotesize
\begin{longtable}{@{}c p{.82\textwidth}@{}}
\caption{Spectres conjoints complets du lecteur \(K_\Omega\) dans les treize
multisections canoniques.}
\label{tab:espectros-komega-multisecciones-rev11}\\
\toprule
\(M\)&\(\operatorname{Spec}_{\rm conj}\mathbf K_M^\Omega\)\\
\midrule
\endfirsthead
\toprule
\(M\)&\(\operatorname{Spec}_{\rm conj}\mathbf K_M^\Omega\)\\
\midrule
\endhead
\(\ell_0\)&\((80,54,6)^{\times1}\)\\
\(\ell_2\)&\((80,54,6)^{\times8}\)\\
\(\ell_4\)&\((24,54,2)^{\times1};(56,54,5)^{\times4};(80,54,6)^{\times11}\)\\
\(\nu_1\)&\((40,54,4)^{\times1};(80,54,6)^{\times1}\)\\
\(\nu_2\)&\((4,54,2)^{\times1};(80,54,6)^{\times3}\)\\
\(\nu_3\)&\((36,54,4)^{\times1};(56,54,5)^{\times1};(80,54,6)^{\times4}\)\\
\(\nu_4\)&\((36,54,4)^{\times1};(56,54,5)^{\times1};(80,54,6)^{\times6}\)\\
\(d_1\)&\((40,54,4)^{\times1};(80,54,6)^{\times1}\)\\
\(d_2\)&\((4,54,2)^{\times1};(80,54,6)^{\times3}\)\\
\(d_3\)&\((36,54,4)^{\times1};(80,54,6)^{\times5}\)\\
\(d_4\)&\((36,54,4)^{\times1};(80,54,6)^{\times7}\)\\
\(u_1\)&\((80,54,6)^{\times4}\)\\
\(u_3\)&\((56,54,5)^{\times3};(80,54,6)^{\times9}\)\\
\bottomrule
\end{longtable}
\endgroup

Les traces peuvent coïncider entre multisections de types différents ; le type,
la multiplicité et le spectre complet demeurent des données de l'opérateur.
Le certificat reproduit ces tableaux ligne par ligne.

\HMTClaim{MD-R-MASS-FIBRE-OPERATORS}
\subsection{Opérateurs de fibre et naturalité du changement de base}

Nous distinguons la fibre canonique \(F_M^{B1}\), formée des positions de
la multisection \(M\) avec leur orientation publiée, et la fibre orientée
\(F_M^{\rm or}\), formée de ses quatre relèvements par position. Pour
\(s\in\{B1,{\rm or}\}\), soit
\(\mathcal H_{M,s}=\bigoplus_{\Gamma\in F_M^s}\mathbb C e_\Gamma\). Nous définissons
\begin{equation}
 B_{M,s}^r e_\Gamma=\kappa_r(\Gamma)e_\Gamma,
 \qquad
 R_{\rm act}^{B_{M,s}^r}
 =\exp\bigl[(\log R_{\rm act})B_{M,s}^r\bigr],
 \qquad r\in\{D,\Omega\}.
 \label{eq:operadores-fibra-two-way-rev11}
\end{equation}
Chaque \(B_{M,s}^r\) est autoadjoint et diagonal dans la base des parcours, de sorte que
\(R_{\rm act}^{B_{M,s}^r}\) est positif. Pour une masse de base positive
\(\widehat M_{M,s}^{(0)}\), la forme opératorielle conditionnée à chaque lecteur est
\begin{equation}
 \widehat M_{M,s}^{\beta,r}
 =R_{\rm act}^{B_{M,s}^r/2}\,
  \widehat M_{M,s}^{(0)}\,
  R_{\rm act}^{B_{M,s}^r/2}.
 \label{eq:ley-operatoria-fibra-rev11}
\end{equation}
Cette expression est positive sans imposer la commutativité. Lorsque
\([\widehat M_{M,s}^{(0)},B_{M,s}^r]=0\), elle se réduit à
\(\widehat M_{M,s}^{(0)}R_{\rm act}^{B_{M,s}^r}\).
Un changement de base unitaire \(U\) transforme
\(B_{M,s}^r\mapsto UB_{M,s}^rU^*\) et conserve spectre, trace, déterminant et
polynôme caractéristique. La fibre électronique canonique \(F_e^{B1}\) est de
rang un et se réduit au scalaire du
\cref{thm:reduccion-electron-two-way-rev11}.
Son relèvement orienté est de rang quatre : les parcours \((++)\) et \((--)\)
partagent le profil électronique, tandis que les orientations mixtes conservent
des profils distincts.

Le déterminant normalisé
\[
 \operatorname{ndet}(R_{\rm act}^{B_{M,s}^r})
 =R_{\rm act}^{\operatorname{Tr}(B_{M,s}^r)/\dim F_M^s}
\]
est naturel par permutation de la fibre. L'additivité logarithmique le distingue
comme moyenne géométrique, mais il n'est pas utilisé comme substitut de la réalisation physique.
Le résultat exact du raffinement reste la paire
\((B_{M,s}^D,B_{M,s}^\Omega)\) ; la section scalaire d'une espèce s'obtient
par le morphisme saveur--génération--parcours déjà construit dans
\cref{eq:morfismo-sabor-ruta-cerrado-rev2,eq:realizacion-fisica-sabor-ruta-cerrada-rev2} :
il sélectionne le parcours persistant par des données non métrologiques, puis applique le
caractère de sa signature. Une fonctionnelle scalaire arbitraire
\(S(B_{M,s}^D,B_{M,s}^\Omega)\) ne définit pas cette sélection et ne demeure qu'un
contrôle empêchant la réintroduction de la masse cible.

\subsection{Échelle absolue, correction relative et sensibilités}

La généalogie de la droite d'action s'écrit
\[
 \Sigma_H=\varphi\alpha_{\rm HMT}^{16},
 \qquad \operatorname{ord}_{10}(\Sigma_H)=-34,
\]
\[
 \delta_{2w}=H_5-\eta_{\rm ret}
 =\frac{90}{\pi}\mathscr C_\pi(\alpha_{\rm HMT}),
 \qquad R_{\rm act}=\frac{H_5}{\eta_{\rm ret}},
\]
\[
 \hbar_{\rm pre}=H_5\,10^{-34}\mathcal U_S,
 \qquad
 \hbar_{\rm ret}=\eta_{\rm ret}\,10^{-34}\mathcal U_S.
\]
Le facteur \(10^{-34}\mathcal U_S\) fixe la section absolue et se simplifie dans les
rapports. Pour le lecteur dodécaphasique,
\begin{equation}
 \log\!\left[
 \frac{(m_X^{\beta,\Omega}/m_e^{\beta})}
 {(m_X^{(0)}/m_e^{(0)})}\right]
 =\bigl[(\Omega_X-80)+(P_{6,X}-6)\Delta_4\bigr]
 \log R_{\rm act}.
 \label{eq:razon-masa-omega-rev11}
\end{equation}
Le terme commun \(-54\alpha_{\rm HMT}\) se simplifie également. Pour le lecteur
chronologique,
\begin{equation}
 \log\!\left[
 \frac{(m_X^{\beta,D}/m_e^{\beta})}
 {(m_X^{(0)}/m_e^{(0)})}\right]
 =\bigl[\kappa_D(\Gamma_X)-\kappa_e\bigr]\log R_{\rm act}.
 \label{eq:razon-masa-D-rev11}
\end{equation}
Avec \(\log R_{\rm act}=6.932817628081925\ldots\times10^{-10}\),
\[
 \frac{\partial\log m}{\partial\Omega}=6.9328\times10^{-10},
 \quad
 \frac{\partial\log m}{\partial P_6}
 =\Delta_4\log R_{\rm act}=7.7090\times10^{-14},
\]
\[
 \frac{\partial\log m}{\partial\alpha}
 =-54\log R_{\rm act}=-3.74372\times10^{-8},
 \quad
 \left.\frac{\partial\log m}{\partial\Delta_4}\right|_e
 =6\log R_{\rm act}=4.15969\times10^{-9}.
\]
Ces dérivées sont des sensibilités du modèle et restent distinctes de
l'incertitude expérimentale d'une masse de comparaison.

\subsection{Appariement physique et secteur nul}

L'appariement typé suit la chaîne
\[
 \begin{aligned}
 &\text{registre physique typé}
 \longrightarrow\text{multisection ou représentation}
 \longrightarrow\text{fibre de routes}\longrightarrow\text{registre}\\
 &\longrightarrow\text{signature et caractère}
 \longrightarrow(B_{M,s}^D,B_{M,s}^\Omega)
 \longrightarrow\text{carte dimensionnelle}
 \longrightarrow\text{confrontation externe}.
 \end{aligned}
\]
L'inventaire physique typé distingue les enregistrements par objet, représentation et codomaine ; les 324 parcours constituent
le moteur interne ; les 471 lignes du catalogue PDG sont des propriétés externes.
Ces cardinaux appartiennent à des domaines différents. La flèche est construite
à partir des descripteurs physiques non métrologiques, de la multisection et du parcours ; une
recherche par proximité avec une masse relève du régime d'étalonnage.

Le secteur nul bifurque avant le caractère positif. Pour le photon et les gluons,
\[
 \kappa_{\rm null}=\mathrm{N/A},
 \qquad m_{\rm null}=0.
\]
L'expression \(R_{\rm act}^{0}=1\) est l'unité du caractère positif. Les
composés concatènent les
enregistrements avant d'évaluer la signature ; les secteurs topologiques conservent fusion et
tresse ; une section virtuelle est classée selon son horizon fini de
prolongement.

\subsection{Indépendance du constructeur de parcours vis-à-vis des masses et clôture de l'opérateur}

La vérification structurelle utilise exclusivement le dénombrement des parcours,
l'enregistrement \(81\times108\) et le transducteur des 81 positions. Elle vérifie leur
couplage ligne par ligne, la couverture
\(81/56/13/324\), les naturalités, le triplet électronique \((80,54,6)\), les
dénombrements (14) et (9), la coïncidence (18/324), les
(40) paires, les descentes différenciées dans \(B1\) et dans les 324 orientations,
la périodicité 54, l'inversion par décalage de 27 et les ablations de feuillet. Le
secteur nul appartient au domaine extérieur au caractère positif. Le calcul est
effectué avant l'introduction de masses, de CODATA, de PDG, de résidus métrologiques ou
d'étiquettes nominales.

Ce certificat n'évalue ni \(\alpha_{\rm HMT}\), ni \(\Delta_4\),
ni \(R_{\rm act}\), ni \(\kappa_e\), ni la carte MeV. Ces évaluations relèvent de
la chaîne arithmétique de la droite d'action et de la réalisation énergétique du
\cref{thm:reduccion-electron-two-way-rev11} ; cette chaîne arithmétique demeure
distincte de la vérification combinatoire des lecteurs.

La comparaison métrologique applique ensuite les opérateurs à l'inventaire typé,
déclare unité, schéma, échelle, date et incertitude, puis effectue des permutations
et une validation par mise en réserve de multisections complètes. Le résultat physique est
régi par le domaine \(81/56/13/324\), par les deux lecteurs de parcours et par
les opérateurs de fibre.

\begin{statusbox}[title={Domaine exact et réalisation physique achevée}]
L'atlas, la chaîne \(C_{81}\to\mathcal L_{108}\) de l'enregistrement chronologique
enrichi de \(108\) itérations,
la droite d'action,
les lecteurs \(K_D,K_\Omega\), leur réduction électronique et les opérateurs de
fibre sont des résultats internes exacts. La paire
\(\mathbf K=(K_D,K_\Omega)\) est le résultat complet sur chaque parcours et
détermine les deux opérateurs diagonaux de toute fibre. Dans la fibre
électronique, les deux lecteurs coïncident. Dans les fibres non centrales, la paire et son
spectre préservent le raffinement complet, tandis que
\(\mathcal S_{\rm phys}\) relie branche, génération, signature fine et, le cas
échéant, orbite de couleur au parcours persistant. Le caractère de ce parcours
donne la coordonnée scalaire avant la confrontation métrologique. Il n'existe donc pas
de disjonction entre opérateur achevé et masse en attente : ce sont deux lectures
typées de la même généalogie. Les deux lecteurs sont codirecteurs, et aucune
proximité avec une masse externe n'intervient dans la sélection.
\end{statusbox}
\HMTClaim{MD-R-MASS-PHYSICAL-SELECTION-CLOSED}
